\section{模型假设}

为突出购电、光伏与储能之间的主要矛盾，并保证四问的信息边界严格一致，本文作如下假设。每条假设都将在相应模型中显式使用。

\begin{enumerate}[label=\arabic*., itemsep=2pt, topsep=3pt]
	\item \textbf{调度周期与时间尺度}：以 $10$ min 为一个调度时段，一天划分为 $T=144$ 个时段，单时段长度 $\Delta t=1/6$ h。附件中标注 $0{:}00{+}1$ 的列即当天 $24{:}00$。
	\item \textbf{不允许向外部电网售电}：题目仅给出微网从外网购电的机制，未给出反向售电的交易规则，因此外网购电量恒为非负，光伏富余只能用于储能充电或弃光。
	\item \textbf{储能建模的简化}：储能充放电效率在全年内恒为 $\eta_c=\eta_d=0.9$（口径敏感性见 \ref{sec:efficiency} 节），不计自放电、容量衰减、循环寿命成本与启停损耗；同一时段不会同时充电与放电（由词典序目标自动排除）。
	\item \textbf{储能功率约束按能量折算}：最大充放电功率 $5000$ kW 对应单时段最大充放电量 $5000\Delta t=833.33$ kWh。
	\item \textbf{初始储电量}：$2025$ 年 $1$ 月 $1$ 日 $0{:}00$ 储电量为 $6000$ kWh；$1$ 月 $31$ 天作为储能状态与预测误差的预热期，正式结果自 $2$ 月 $1$ 日起输出，这既与官方结果模板的日期范围一致，也避免冷启动误差污染统计。
	\item \textbf{负载预测口径}：题目只为光伏提供了专门预报数据，未给出负载预报。因此本文用历史实际负载构造因果预测，方法族为持续法、$7/14/30$ 日滚动均值、$14/30$ 日线性加权均值与同星期均值，并在历史执行块上滚动在线选型，选型标准为平均绝对误差。
	\item \textbf{附件 1 作为冷启动先验}：$1$ 月 $1$ 日之前无任何历史观测，此时以附件 1 给出的典型日负载与光伏预测曲线作为先验基预测。该先验仅在历史观测为空的当天使用，一旦有了至少一天观测便自动失效。
	\item \textbf{实时储能反馈}：计划锁定后，日内允许储能依据\emph{当前时刻已经发生的}真实光伏与当前储电量做 $10$ min 级调整，调整顺序为“先撤充电、再增放电、最后紧急购电”，不得借助任何未来信息。
	\item \textbf{终端储电量策略}：题目仅对问题一明确要求 $S_{0{:}00}=S_{24{:}00}$。本文主模型在问题二、三的计划层同样采用“计划终端储电量等于当日实际初始储电量”的日循环口径（记为 daily\_cycle），以保证储能长期可持续运行；同时给出不做终端约束（free）的敏感性结果，二者的差异在 \ref{sec:terminal} 节报告。
	\item \textbf{问题四的价格信息集}：价格完全已知情形视为理想化基准；因果情形的价格预测只使用决策时刻之前已经发生的附件 4 真实电价，跨日部分由历史同刻价格外推。
\end{enumerate}

\section{符号说明}

本文使用的主要符号及其含义如表~\ref{tab:symbols} 所示。表中未列出的局部符号将在正文首次出现处说明。

\begin{table}[htbp]
	\caption{主要符号说明}
	\label{tab:symbols}
	\centering
	\zihao{-5}
	\begin{tabularx}{0.94\textwidth}{l l X}
		\toprule
		符号 & 单位 & 含义 \\
		\midrule
		$d,\ t$ & --- & 日期下标与日内 $10$ min 时段下标（$t=1,\dots,144$）\\
		$\Delta t$ & h & 单个调度时段长度，$\Delta t=1/6$ \\
		$L_{d,t}$ & kW & 小区负载功率 \\
		$P_{d,t}$ & kW & 光伏实际发电功率 \\
		$\hat P_{d,t}$ & kW & 决策时刻可用的光伏预测功率 \\
		$\pi_{d,t}$ & 元/kWh & 外网正常购电单价 \\
		$G^{p}_{d,t}$ & kWh & $0{:}00$ 制定的计划购电量（第一阶段变量）\\
		$G^{a}_{d,t}$ & kWh & 执行前最后一次有效调整后的购电量（第二阶段变量）\\
		$C_{d,t}$ & kWh & 母线侧进入储能的充电量 \\
		$H_{d,t}$ & kWh & 储能送往母线的放电量 \\
		$S_{d,t}$ & kWh & 第 $t$ 时段结束时储能内部储电量（状态变量）\\
		$W_{d,t}$ & kWh & 未被利用的弃光电量 \\
		$E_{d,t}$ & kWh & 供电不足时的紧急购电量 \\
		$u_{d,t},\ v_{d,t}$ & kWh & 调整购电量相对计划的上调量 $(G^a-G^p)^+$ 与下调量 $(G^p-G^a)^+$ \\
		$\eta_c,\ \eta_d$ & --- & 充电效率与放电效率，均取 $0.9$ \\
		$S_{\min},\ S_{\max}$ & kWh & 储电量上下限，$1200$ 与 $10800$ \\
		$\bar E$ & kWh & 单时段充放电量上限，$\bar E=5000\Delta t=833.33$ \\
		$\mathcal{I}_d$ & --- & 第 $d$ 天决策时刻的信息集 \\
		\bottomrule
	\end{tabularx}
\end{table}
